\chapter*{Scope and reading guide}
\addcontentsline{toc}{chapter}{Scope and reading guide}
\markboth{Scope and reading guide}{}

This supplement develops symmetric powers over arbitrary commutative
rings and Witt decomposition and cancellation over fields of
characteristic different from two. The route through ALG-01 below
points to the existing proofs and identifies where the new chapters
enter.

{\setlength{\fboxsep}{0.8em}\setlength{\fboxrule}{0.35pt}%
\noindent\fcolorbox{rulegreen}{palegreen}{%
\begin{minipage}{\dimexpr\linewidth-2\fboxsep-2\fboxrule\relax}
\textbf{Result keys.} \texttt{LA01-C05.01} belongs to module
\texttt{ALG-01-C05}. The source map records its conventions,
dependencies, and PDF and TeX locations. Keys beginning with
\texttt{K05} refer to result groups in the released D05 notes.
\end{minipage}}\par}

\section*{The six core modules}

\paragraph{C01 \enspace Vector spaces, modules, and dimension.}
D05, groups K05.02--03 and K05.06, proves the field results on bases,
rank--nullity, and quotients \cite{d05}.
\emph{Algebra Core} supplies module quotients, exactness, and split
sequences \cite[Ch.~1]{p103}. D07 treats modules over a PID \cite[Ch.~16]{d07}.

\paragraph{C02 \enspace Duality and forms.}
D05, groups K05.07 and K05.11, gives dual bases, the dual of a quotient,
scalar conventions on forms, and nondegeneracy.
Group K05.13 gives the adjoint on finite-dimensional real and complex
inner-product spaces.

\paragraph{C03 \enspace Determinants and canonical forms.}
D05, groups K05.05, K05.08, and K05.10, gives determinants, Cayley--Hamilton,
primary decomposition, and rational and Jordan forms.
Jordan form uses a splitting field. For the Smith form over a PID,
see D07, Chapter~16.

\paragraph{C04 \enspace Spectral and singular-value methods.}
D05, groups K05.12--14, proves the finite-dimensional real and complex
spectral, SVD, and polar results. For the Hilbert-space partial-isometry
version see \cite{d16}, and for the min--max and matrix perturbation
bounds see \cite{p206}.

\paragraph{C05 \enspace Tensor, exterior, and symmetric powers.}
The balanced tensor construction is in \cite[Ch.~1]{p103}. D05,
group K05.15, treats the field case. The comparison
\(S\otimes_R\bigwedge_R^n M\simeq\bigwedge_S^n(S\otimes_R M)\)
is the degree-\(n\) comparison in \cite[RK-03 and RK-04]{commalg}.
Chapter~\ref{ch:la01-symmetric} develops the symmetric powers over rings.

\paragraph{C06 \enspace Quadratic forms and Witt theory.}
D05, group K05.11, provides polarization, diagonalization, and the
orthogonal-complement results. Chapter~\ref{ch:la01-witt} gives
hyperbolic splitting and Witt cancellation with the characteristic
hypothesis and the nondegenerate scope planned for NT-06.
For the classical treatment see \cite[Ch.~I, Sec.~4]{lam-witt}.
